\documentclass[10pt,letterpaper]{article}

% --- CONTROL KNOBS ---
\def\DocumentMargin{0.5in}
\def\MainColorRGB{0,0,255}
\def\MainFontName{IBMPlexSans}
% ---------------------

\usepackage[margin=\DocumentMargin]{geometry}
\usepackage{fontspec}
\usepackage{xcolor}
\usepackage{tabularx}
\usepackage{enumitem}
\usepackage{ragged2e}
\usepackage{microtype}
\usepackage{array}
\usepackage{hyperref}

\definecolor{PureBlue}{RGB}{\MainColorRGB}

\setmainfont{\MainFontName}[
  Extension=.otf,
  UprightFont=*-Regular,
  BoldFont=*-SemiBold,
  ItalicFont=*-Italic,
  BoldItalicFont=*-SemiBoldItalic
]

\urlstyle{same}
\hypersetup{
  colorlinks=true,
  urlcolor=PureBlue,
  linkcolor=PureBlue,
  citecolor=PureBlue,
  pdftitle={Onri Jay Benally Quantum Hardware Systems Resume},
  pdfauthor={Onri Jay Benally}
}

\pagestyle{empty}
\setlength{\parindent}{0pt}
\setlength{\parskip}{2pt}
\RaggedRight
\microtypesetup{protrusion=true, expansion=false}

\newcommand{\name}[1]{%
  {\fontsize{24}{26}\selectfont\bfseries\color{PureBlue} #1}\par
}

\newcommand{\roleline}[1]{%
  {\fontsize{10.5}{12}\selectfont\bfseries #1}\par
}

\newcommand{\sectionheading}[1]{%
  \vspace{4pt}%
  {\fontsize{9.5}{10.5}\selectfont\bfseries\color{PureBlue}\MakeUppercase{#1}}\par
  \vspace{-6pt}{\color{PureBlue}\rule{\linewidth}{0.45pt}}\vspace{1pt}
}

\newcommand{\entry}[4]{%
  \textbf{#1}\hfill {\color{PureBlue}#2}\par
  {\fontsize{8.7}{10}\selectfont #3}\par
  \vspace{-2pt}#4\vspace{1pt}
}

\newenvironment{compactitems}{%
  \begin{itemize}[leftmargin=12pt,itemsep=1.0pt,topsep=1pt,parsep=0pt,partopsep=0pt]
}{%
  \end{itemize}
}

\newcolumntype{Y}{>{\RaggedRight\arraybackslash}X}

\begin{document}

\name{ONRI JAY BENALLY}
\roleline{Quantum Hardware Systems Engineer \textmd{| Magnetics | Frontier Hardware Knowledge Repositories | Quantum Chip-Patterning Workflows}}
{\fontsize{8.7}{10}\selectfont
Minneapolis, MN \ | \ \href{mailto:ojbs.name@gmail.com}{ojbs.name@gmail.com} \ | \ \href{tel:+16514343799}{(651) 434-3799} \ | \ \href{https://tinyurl.com/OJB-LinkedIn}{LinkedIn} \ | \ \href{https://github.com/OJB-Quantum}{GitHub} \ | \ \href{https://z.umn.edu/OJB-Google-Scholar}{Google Scholar} \ | \ \href{https://orcid.org/0000-0002-8391-9105}{ORCID}
}

\sectionheading{Strategic Profile}
\fontsize{8.9}{10.15}\selectfont
Build and de-risk the device platforms, cryogenic test infrastructure, chip-patterning workflows, and technical knowledge infrastructure needed to move quantum hardware from
component-level research toward scalable system architectures.

8+ years of hands-on experience across nanofabrication, cryogenic hardware testing, spintronic magnetics, superconducting resonator metrology, quantum chip patterning, and executive-facing technical roadmaps for quantum stakeholders.

\sectionheading{Core Engineering Value}
\begin{tabularx}{\linewidth}{@{}p{0.31\linewidth}Y@{}}
\textbf{Scalable quantum systems} & Cryogenic MRAM, hybrid QPU hardware, magnetic cluster qubit prototypes, superconducting resonator chiplets, quantum-limited-amplifier test conditions, and dilution-fridge workflows. \\
\textbf{Magnetics-to-quantum hardware bridge} & MTJ and p-MTJ arrays, SOT-MRAM, cryogenic SOT CRAM, CoFeB-based quantum magnonic buses, magnetic nanocluster deposition/patterning, and spin-orbit torque (SOT) material device development. \\
\textbf{Frontier hardware knowledge systems} & Reusable documentation, GitHub/Markdown tutorials, photomask layout training, scientific-computing examples, Blender 3D quantum hardware illustrations, and repository-style chip-patterning guidance. \\
\textbf{Technical strategy and translation} & Trusted technical contact for senior sell-side equity analysts and quantum industry leaders; invited lecturer/keynote speaker for quantum hardware, scientific computing, and quantum education audiences. \\
\end{tabularx}

\sectionheading{Technical Systems}
\begin{tabularx}{\linewidth}{@{}p{0.2\linewidth}Y@{}}
\textbf{Cryogenic and quantum\newline hardware} & Bluefors XLDsl dilution refrigerator, mK testing, Nb thin films, CPW resonators, superconducting resonator chiplets, grAl/MoGe high-kinetic-inductance resonators, UHV in-situ shadow-mask patterning, packaging/wiring, low-temperature metrology. \\
\textbf{Nanofabrication\newline and metrology} & Raith EBPG-5000+ 100 kV, Karl Suss MA6, Heidelberg DWL200, PECVD, ALD, RIE, IBE, sputtering, electron-beam evaporation, SEM, AFM, profilometry, LASER scanning microscopy, metal hardmasks, MTJ process flows. \\
\textbf{Scientific computing\newline and repositories} & Python, Julia, Go, CUDA Python/CuPy, Qiskit, Quantum Metal, MuMax3, Elmer FEM, Gmsh, ParaView, KLayout, LinkCAD, Genisys BEAMER, Git/GitHub, Markdown, LaTeX, Linux, IBM Granite, Google Colab, Ollama, SolidWorks. \\
\end{tabularx}

\sectionheading{Experience}
\entry{NSF Graduate Research Fellow}{Aug 2024 - Present}{University of Minnesota - Nano Magnetism and Quantum Spintronics Lab, ECE Dept. \ | \ Minneapolis, MN}{%
\begin{compactitems}
  \item Engineering cryogenic MRAM and hybrid QPU hardware integrating advanced spintronics materials for ultra-low-Joule qubit-control and readout-sampling concepts across mK-to-K operating regimes.
  \item Collaborate with Idaho National Laboratory, MIT Lincoln Laboratory, and the UMN Gang Qiu Group on a magnetic cluster qubit platform; responsible for quantum device design prototypes and qubit-device patterning recipes.
  \item Developed experiment plans and device recipes for Ni4W cryogenic SOT devices, cryogenic SOT CRAM, CoFeB-based quantum magnonic buses, 30 nm MTJs with metal hardmasks, grAl/MoGe resonators, and Tb magnetic nanocluster deposition/patterning.
  \item Developed a novel UHV-compatible in-situ metal shadow-mask patterning system for high-throughput device patterning and testing within a modular system rated at $10^{-9}$ Torr.
  \item Serve as Principal Investigator for an ASU Quantum Collaborative-funded quantum hardware and quantum science education project in the Navajo language.
  \item Act as a confidential technical roadmap contact for senior sell-side equity analysts and quantum industry leaders; deliver faculty-level talks on quantum-limited amplifiers, cryogenic memories, control systems, cryogenic testing setups, and IBM-based scientific computing for magnetics.
\end{compactitems}}

\entry{Graduate Research Assistant}{Jun 2022 - Present}{University of Minnesota - Nano Magnetism and Quantum Spintronics Lab, ECE Dept. \ | \ Minneapolis, MN}{%
\begin{compactitems}
  \item Recognized as resident nanodevice patterning expert for robust, solid-state, integrated spintronic chips; led high-quality superconducting MTJ arrays and hybrid magnetic qubit work across TaN, IrO$_{2}$, RuO$_{2}$, and Nb materials with device sizes from 80 nm to 6 µm.
  \item Trained and mentored academic researchers in high-yield quantum device fabrication, electron microscopy, and photomask design/layout; supplemented training with \href{https://github.com/OJB-Quantum}{GitHub repository} Markdown tutorials.
  \item Delivered invited internal and external lecture/workshop series on scientific illustration aided by micrographs, real quantum data, and scientific-computing results; created Colab workflow examples for Quantum Metal, formerly Qiskit Metal, enabling reproducible superconducting and hybrid quantum device layout tutorials for quantum hardware learners.
\end{compactitems}}

\entry{Quantum Hardware Engineer - PhD Intern}{May 2024 - Aug 2024}{IBM Thomas J. Watson Research Center - Kevin Tien Group \ | \ Yorktown Heights, NY}{%
\begin{compactitems}
  \item Automated a high-throughput in-situ low-temperature metrology process for high quality factor superconducting Nb thin films and CPW resonators supporting quantum-centric supercomputing systems.
  \item Developed 10 mK test setups for superconducting resonator chiplets from experimental fabrication splits to mimic Josephson Traveling Wave Parametric Amplifier working conditions for integration.
  \item Sole hardware-engineer intern permitted to assemble and operate a Bluefors XLDsl dilution refrigerator hands-on for experiments and configure IBM data servers with Linux.
  \item Produced extensive experimental-design and setup documentation; recognized for high-quality presentations, hands-on quantum measurement lab work, and quantum hardware education/illustration initiatives.
\end{compactitems}}

\sectionheading{Selected Earlier Hardware Work}
\entry{Research Assistant - UMN Nano Magnetism and Quantum Spintronics Lab}{2020 - 2022}{Minneapolis, MN}{%
\begin{compactitems}
  \item Trained 11 PhD-level academic users and 1 faculty member on high-yield quantum device fabrication and electron microscopy as a highly skilled undergraduate researcher; developed patterning protocols for p-MTJ arrays in SOT-MRAM, CRAM, and magnetic brain-interface chips using PtSn, PtSi, CoPd, CoSi, TaN, FePd, and CoFe alloys.
\end{compactitems}}

\entry{Research Assistant - REU - UMN Quantum Devices and Materials Lab}{2017}{Minneapolis, MN}{%
\begin{compactitems}
  \item Optimized process development for III-V InSb Majorana nanowire devices and 100 nm spin-valves for Majorana quasiparticle studies and topological quantum logic-gate research.
\end{compactitems}}

\entry{Research Assistant - REU - UMN Nano Magnetism and Quantum Spintronics Lab}{2018}{Minneapolis, MN}{%
\begin{compactitems}
  \item Developed process recipes for microscale MTJ arrays and Hall bars; obtained up to 50\% magnetoresistance ratio in FePd and Co/Pd-based MTJs.
\end{compactitems}}

\entry{Lead Volunteer - Electric-Vehicle and Power Hardware}{2016 - 2017}{}{%
\begin{compactitems}
  \item Inspected/tested 60+ kW EV subsystems and led field-demonstrated electric ATV, battery, powertrain, robotics, PV-array, and modular solar-charging hardware projects.
\end{compactitems}}

\sectionheading{Education and Credentials}
\begin{tabularx}{\linewidth}{@{}p{0.28\linewidth}Y@{}}
\textbf{Ph.D., Electrical Engineering} & University of Minnesota-Twin Cities, quantum hardware and spintronics, expected Dec 2027. \\
\textbf{B.S.Md.S., Multidisciplinary Studies} & University of Minnesota-Twin Cities, Dec 2021. \\
\textbf{A.S., General Studies, Physics} & Utah State University-Eastern, Dec 2016. \\
\textbf{Selected certifications} & IBM Certified Quantum Developer - Qiskit v0.2x; IBM Quantum Business Foundations; Code Generation and Optimization Using IBM Granite; IBM Growth Behaviors. Full training portfolio: \href{https://z.umn.edu/OJB-Certifications}{z.umn.edu/OJB-Certifications}. Credential profile: \href{https://www.credly.com/users/onri-jay-benally}{Credly}. \\
\end{tabularx}

\sectionheading{Leadership, Strategy, and Communication}
\begin{tabularx}{\linewidth}{@{}p{0.28\linewidth}Y@{}}
\textbf{ASU Quantum Collaborative} & 2025 Seed Fund Project Principal Investigator; 2024 Steering Committee Member. \\
\textbf{Technical community leadership} & 2026 IEEE Magnetics Society Twin Cities Chapter Vice President; Quantum Device Consortium Quantum Metal community advocate/workflow maintainer; Open Source Hardware Association member; QuantumGrad mentor and lecturer. \\
\textbf{Selected invited/keynote speaking} & Quantum Device Consortium speaker, 2026; Harvard SEAS invited lecturer/keynote speaker, 2026; University of Maryland Q$^{3}$ Joint Quantum Initiative QREATE keynote speaker, 2025; Quantum Collaborative Summit keynote speaker, 2023; UMN Quantum + Chips invited lecturer, 2023. \\
\textbf{Education and outreach} & Trained more than 25 academic researchers in cleanroom nanodevice fabrication; created quantum and hardware career-readiness content, STEM demonstrations, and public-facing quantum hardware models. \\
\end{tabularx}

\sectionheading{Selected Publications and Technical Output}
\begin{tabularx}{\linewidth}{@{}p{0.28\linewidth}Y@{}}
\textbf{Publication record} & 17 papers/abstracts spanning MRAM, MTJs, SOT materials, voltage-controlled exchange coupling, superconducting/spintronic hardware, and embedded quantum computing architecture. Full publication profile: \href{https://z.umn.edu/OJB-Google-Scholar}{Google Scholar}. \\
\textbf{Main-author architecture} & ``Magnetic random-access memory (MRAM) for embedded quantum computing hardware,'' GOMACTech submitted abstract, 2023. \\
\textbf{Selected co-authored venues} & Nature Communications, Advanced Materials, Advanced Electronic Materials, Applied Physics Letters, Physical Review B, Physical Review Materials, Nano Letters, Scientific Reports, EIPBN. \\
\textbf{Public technical footprints} & \href{https://github.com/OJB-Quantum}{GitHub/OJB-Quantum}, \href{https://z.umn.edu/OJB-Google-Scholar}{Google Scholar}, \href{https://orcid.org/0000-0002-8391-9105}{ORCID}, \href{https://tinyurl.com/OJB-Projects-001}{hardware projects portfolio}. \\
\end{tabularx}

\sectionheading{Selected Recognition}
\fontsize{8.9}{10.2}\selectfont
\textbf{Advanced Materials journal cover feature (2025)} \ | \ ECE Department media recognition (2025) \ | \textbf{\ NSF Graduate Research Fellowship (2024)} \ | \ IBM media and executive recognition for quantum hardware illustration/education (2024) \ | \ Womanium Global Quantum Scholarship (2023) \ | \ Raith GmbH media recognition (2021) \ | \ Thank-a-Teacher Award (2021) \ | \ IBM Quantum Qubit by Qubit full sponsorship (2020) \ | \ NSF MRSEC Fellowships - REU (2017, 2018) \ | \ Chief Manuelito Scholar Award and Presidential Scholar Award (2014-2017)

\end{document}
